% !PRINTED 307
\weberpage{307}

%<<CONT
$\varrho$ runs in (7) and (8) through either the sequence of the quadratic residues or that of the non-residues, so that for $p = 7$ or 11 one obtains \gesperrt{in each case two groups of index $p$.}

That the systems of substitutions collected in (3), (7), (8) really constitute groups can be shown by direct composition in various ways. One reaches this proof more simply, however, by forming functions of the $v_z$, § 78, (7), which remain unchanged under the substitutions of these systems, and which under the substitutions of $\mathfrak{L}_0$ take only $p$ different values in all.

The systems (3), (7), (8) can, as follows from their mode of origin, be composed by repeated application of $B, C, U, U'$, if $U, U'$ are any two special substitutions $W, W'$. For $p = 5$ the substitution $U^2$, and for $p = 7, 11$ the substitution $UB$, belongs to the $W'$, so that $U'$ can then be omitted. If we choose $U$ quite arbitrarily, and take for the primitive root $g$ in $C$, for $p = 5, 7, 11$ respectively, $2, -2, 2$, then we can compose the group (3) from:
\[
B = \begin{pmatrix} 0, & 1 \\ -1, & 0 \end{pmatrix}, \quad C = \begin{pmatrix} 2, & 0 \\ 0, & 3 \end{pmatrix}, \quad U = \begin{pmatrix} 1, & 1 \\ 2, & 3 \end{pmatrix} \quad (p = 5),
\]
the group (7) from:
\[
B = \begin{pmatrix} 0, & 1 \\ -1, & 0 \end{pmatrix}, \quad C = \begin{pmatrix} 2, & 0 \\ 0, & 4 \end{pmatrix}, \quad U = \begin{pmatrix} 1, & \pm 1 \\ \pm 2, & 3 \end{pmatrix} \quad (p = 7),
\]
the group (8) from:
\[
B = \begin{pmatrix} 0, & 1 \\ -1, & 0 \end{pmatrix}, \quad C = \begin{pmatrix} 2, & 0 \\ 0, & 6 \end{pmatrix}, \quad U = \begin{pmatrix} 1, & \pm 1 \\ \pm 1, & 2 \end{pmatrix} \quad (p = 11),
\]
where in the two last cases the two different groups mentioned above arise from the double signs.

We now collect the permutations of the indices $z$ that correspond to these substitutions [§ 79, (3)].

$p = 5$:
\[
\begin{array}{ccccccl}
z = \infty, & 0, & 1, & 2, & 3, & 4 & \\
0, & \infty, & 4, & 2, & 3, & 1 & (B)\\
\infty, & 0, & 4, & 3, & 2, & 1 & (C)\\
4, & 1, & 2, & 0, & \infty, & 3 & (U).
\end{array}
\]
Thus by $B, C, U$ the index pairs
\[
(\infty, 0), \quad (1, 4), \quad (2, 3)
\]
%CONT>>
